\section{Introduction}
\label{sec:intro}

Few-shot prompting is the default way to turn a large language model (LLM) into a classifier: paste a handful of labelled examples into the prompt, then ask for the label of a new input~\citep{brown2020language}. The demonstrations are usually picked by hand. As a result, they can carry statistics that have nothing to do with the task. Most demonstrations may share one label. The last few may share a label. One label may come with longer inputs than the others. If the model picks up these statistics, a deployed classifier can quietly change its behaviour depending on which examples someone happened to paste.

Early evidence suggested this is a serious problem. \citet{zhao2021calibrate} showed that GPT-3 is pulled toward the majority label and toward the label of the last demonstrations, and that this makes few-shot accuracy unstable. \citet{lu2022fantastically} showed that the order of demonstrations alone can move accuracy from near chance to near state of the art. \citet{fei2023mitigating} identified a further domain-label bias. These findings motivated a family of calibration methods~\citep{zhao2021calibrate,fei2023mitigating,zhou2024batch}.

{\bf Do these biases still hold?} Most of this evidence comes from GPT-3-era base models, and most studies vary one factor at a time. Practitioners now mostly use instruction-tuned models~\citep{ouyang2022training} from 1B to tens of billions of parameters. For these models we lack a side-by-side answer to two simple questions. First, which spurious demonstration statistics matter, and by how much? Second, does the answer change with model scale? A common intuition, which we turn into explicit hypotheses in \secref{sec:method}, is that frequency skew causes the largest bias (15--25 percentage points), recency comes second (8--15 points), label--length correlation causes a moderate bias that scale cannot fix, larger models correct skew but amplify recency, and calibration gets worse under every bias.

{\bf Our approach.} We test these claims with one paired protocol. We apply four manipulations to a 10-shot demonstration set: frequency skew, recency, primacy (label position), and label--length correlation (\tabref{tab:conditions}). We compare each against a balanced \clean set on \sst, \rte, \agnews and \mrpc. The label favoured by the manipulation, the \emph{target} $T$, is rotated through every label. A model's intrinsic preference for one label therefore cancels, and the shift toward $T$ measures the bias directly. Every model sees the same 10,000 prompts. We score labels by their exact next-token probabilities, which removes parsing errors and gives calibrated-probability metrics for free. We evaluate eight local models: a dense \qwen ladder (0.5B, 1.5B, 3B, 7B, 14B), a Qwen2.5-1.5B base model, \llamaone-Instruct and \llamaeight-Instruct.

{\bf What we find.} Every bias is far smaller than the hypotheses predict (\figref{fig:shift}). Frequency skew is the largest bias in small models, at +4.3 to +4.9 points for small \qwen models, +8.9 for the 1.5B base model and +16.9 for \llamaone. It shrinks by 3.7 points per tenfold increase in parameters and is about zero at 14B ($-0.6$) and 8B (+0.6). Recency never produces the predicted pull. Instead, models of 3B and above are slightly \emph{less} likely to predict the label of the final demonstrations ($-1.7$ to $-3.6$ points, Holm-corrected $p<0.05$). Label--length correlation has no reliable effect at any size ($|\shift|\le3.1$ points), and only 1 of 32 calibration-error differences has a confidence interval that excludes zero. Large shifts need extreme manipulations: when all ten demonstrations share one label, models of 1.5B and below shift by 16--36 points, but Qwen2.5-7B shifts by only 3.2. For every model of 3B and above, simply re-drawing a balanced demonstration set flips 4--8\% of predictions. That is more than any systematic bias we measure in these models.

Our contributions are:
\begin{itemize}[leftmargin=*,itemsep=0pt,topsep=0pt]
    \item We build a paired, label-rotated benchmark that measures frequency, recency, primacy and label--length biases side by side with one protocol. Every model sees the same prompts, so all comparisons are paired at the prompt level.
    \item We conduct a scale study over eight models and 80,000 scored prompts. It shows that frequency skew fades with scale and instruction tuning, and that recency reverses sign in models of 3B and above.
    \item We complement the main study with dose--response, length-interaction, calibration and resampling-noise analyses. They locate the remaining bias in small models, extreme manipulations and ambiguous tasks.
    \item We give practical guidance: for current instruction-tuned models of 3B and above, averaging over demonstration sets matters more than controlling label order or length.
\end{itemize}

\Secref{sec:related} reviews related work. \Secref{sec:method} describes the benchmark and metrics. \Secref{sec:results} reports results, \secref{sec:discussion} discusses them along with limitations, and \secref{sec:conclusion} concludes.
